\documentclass[12pt,a4paper]{article}

% ─── Packages ────────────────────────────────────────────────
\usepackage{amsmath,amssymb,amsthm}
\usepackage{graphicx}
\usepackage{booktabs}
\usepackage{natbib}
\usepackage[margin=1in]{geometry}
\usepackage{setspace}
\usepackage{hyperref}
\usepackage{xcolor}
\usepackage{float}
\usepackage{fontspec}           % XeLaTeX for Unicode/CJK support
\usepackage{xeCJK}              % CJK support
\setCJKmainfont{PingFang TC}    % macOS built-in Traditional Chinese font
\usepackage{multirow}
\usepackage{threeparttable}
\usepackage{caption}

\hypersetup{colorlinks=true, linkcolor=blue, citecolor=blue, urlcolor=blue}
\doublespacing

% ─── Theorem environments ────────────────────────────────────
\newtheorem{proposition}{Proposition}
\newtheorem{hypothesis}{Hypothesis}
\newtheorem{definition}{Definition}

% ─── Title ───────────────────────────────────────────────────
\title{Volatility Absorption: The Diminishing Marginal Impact of Market Fear\thanks{We thank seminar participants for helpful comments. Computational support was provided by the VolPred Research System. All errors remain our own. Data and replication code available upon request.}}

\author{Yi-Hao Lai\thanks{Department of Finance, Da-Yeh University. Corresponding author. Email: yhlai@mail.dyu.edu.tw} \and VolPred Research System}

\date{March 2026}

\begin{document}

\maketitle

% ─── Abstract ────────────────────────────────────────────────
\begin{abstract}
\noindent
We document a novel empirical regularity: the marginal impact of fear shocks on realized asset returns diminishes as the ambient fear level rises---a phenomenon we term \textit{volatility absorption}. Using daily data on four assets (SPY, GLD, TLT, 0050.TW) and VIX from 2006 to 2026, we show that the ratio of shocked-day to normal-day absolute returns, normalized by the prevailing VIX level, declines monotonically from 3.16$\times$ in calm markets to 2.32$\times$ in high-fear regimes. This absorption effect is confirmed across three of four assets, with the exception being the Taiwanese equity ETF, consistent with VIX being a US-specific fear gauge. We further decompose shocks by type and find that \textit{endogenous} risks---rate shocks and risk-off flights already reflected in VIX---are heavily absorbed (coefficient $+0.019$), whereas \textit{exogenous} risks such as geopolitical shocks are not ($-0.003$). Event-level evidence from nonfarm payroll releases confirms this pattern: NFP-day volatility is 1.17$\times$ normal in calm markets but indistinguishable from normal days when VIX exceeds 25. These findings have direct implications for hedging, portfolio rebalancing, and volatility-targeting strategy design.

\medskip
\noindent
\textbf{Keywords:} Volatility absorption, VIX, variance risk premium, fear shocks, endogenous risk, exogenous risk, volatility targeting, hedging

\medskip
\noindent
\textbf{JEL Classification:} G11, G12, G14, G17
\end{abstract}

\newpage
\tableofcontents
\newpage

% ══════════════════════════════════════════════════════════════
% SECTION 1: INTRODUCTION
% ══════════════════════════════════════════════════════════════
\section{Introduction}
\label{sec:intro}

Financial markets exhibit a puzzling pattern during sustained crises: after an initial spike in volatility, subsequent shocks of comparable magnitude produce progressively smaller market reactions. During the 2008 Global Financial Crisis, for example, a two-point VIX increase when VIX was already at 60 moved SPY far less, in volatility-adjusted terms, than the same two-point increase when VIX was at 15. Market participants, policymakers, and risk managers have long observed this ``numbness'' informally, yet the phenomenon has received surprisingly little systematic empirical attention.

We formalize and quantify this observation under the label \textit{volatility absorption}: the diminishing marginal impact of fear shocks on realized asset price movements as the ambient level of market fear rises. Our central measure is the \textit{normalized shock impact}---the absolute return on a shock day divided by the prevailing VIX level---and we show that this quantity declines monotonically across VIX regimes. When VIX is below 15 (``calm''), a large VIX shock ($|\Delta\text{VIX}| > 2$) produces an absolute SPY return that is 3.16 times the normal-day average. When VIX is above 25 (``high fear''), the same shock produces only 2.32 times the normal-day average. The difference is economically large and statistically significant.

The absorption effect is not merely a mechanical consequence of higher baseline volatility. After normalizing returns by VIX level, we estimate a cross-sectional regression of normalized shock impact on VIX and find a slope of $-0.00028$, confirming that each additional VIX point reduces the marginal impact of the next shock. This ``paralysis'' coefficient is robust to alternative shock thresholds and sub-period splits.

Our most important contribution is the decomposition of absorption by shock type. We classify shocks into three categories based on the joint behavior of SPY, TLT, and GLD on shock days: (i) \textit{rate shocks} (SPY$\downarrow$, TLT$\downarrow$), in which both equities and bonds fall, typically driven by monetary policy surprises; (ii) \textit{risk-off flights} (SPY$\downarrow$, TLT$\uparrow$), the classic flight-to-quality pattern; and (iii) \textit{geopolitical shocks} (SPY$\downarrow$, GLD$\uparrow$), in which gold rallies as a safe haven while equities decline. The key finding is that endogenous risks---those already reflected in the VIX level, such as rate shocks and risk-off flights---are heavily absorbed (absorption coefficients of $+0.019$ and $+0.007$, respectively), whereas exogenous risks such as geopolitical shocks are \textit{not} absorbed (coefficient $-0.003$). This distinction between endogenous and exogenous absorption has, to our knowledge, not been documented in the prior literature.

Event-level evidence corroborates the aggregate findings. We examine nonfarm payroll (NFP) release days---a scheduled macroeconomic event whose impact should, in theory, depend on the surprise component rather than the ambient fear level. In calm markets (VIX $< 20$), NFP days produce absolute returns 1.17 times normal ($p = 0.037$). But when VIX exceeds 25, the NFP effect vanishes entirely: the shock-to-normal ratio drops to 0.95$\times$ and is statistically indistinguishable from unity. Scheduled macro events, which are endogenous to the information set already priced by VIX, are absorbed in the same way as general VIX shocks.

We extend our analysis to multiple assets and find that the absorption effect is universal across US-traded assets. SPY (slope $-0.00028$), GLD ($-0.00043$), and TLT ($-0.00044$) all exhibit negative normalized regression slopes, confirming diminishing marginal fear impact. The sole exception is the Taiwan equity ETF 0050.TW (slope $+0.00019$), which we interpret as evidence that VIX is a US-specific fear gauge with limited direct relevance for non-US equity markets. Taiwanese equities respond to VIX indirectly through US lead-lag effects, and a US-specific VIX shock may not capture the incremental fear relevant to the Taiwan market.

Our findings connect to several strands of the literature. First, the volatility clustering literature \citep{engle1982, bollerslev1986} documents that high-volatility periods persist, but does not address the diminishing \textit{marginal} impact within those periods. Second, the variance risk premium (VRP) literature \citep{bollerslev2009, carr2009} shows that implied volatility systematically exceeds realized volatility. We find that the VRP narrows at high VIX ($+2.8\%$ annualized versus $+3.5\%$ at low VIX) but remains strictly positive---there is no VRP sign flip, contrary to some prior claims. Third, the literature on the leverage effect \citep{black1976, christie1982} documents asymmetric volatility responses to returns but does not examine how the magnitude of this response varies with the ambient fear level. Fourth, the behavioral finance literature on ``panic fatigue'' or ``crisis habituation'' \citep{kahneman1979, barberis2001} provides a psychological foundation for our empirical findings.

The economic implications of volatility absorption are substantial. Hedging cost-benefit ratios decline from 13.7$\times$ in calm markets to 3.6$\times$ in high-fear regimes, suggesting that the marginal value of incremental hedging is dramatically lower during crises. Portfolio rebalancing adds approximately zero marginal value across VIX regimes. Crucially, the widely-used volatility targeting (VT) strategy with a $12/\text{VIX}$ equity weight naturally accounts for absorption: because exposure is inversely proportional to VIX, the strategy automatically reduces its sensitivity to incremental shocks as fear rises. We show that adding overlays (momentum filters, regime switches) to VT strategies at high VIX levels is not only unnecessary but actively harmful, consistent with the absorption framework.

The remainder of the paper proceeds as follows. Section~\ref{sec:lit} reviews the related literature. Section~\ref{sec:method} describes our methodology. Section~\ref{sec:data} presents the data. Section~\ref{sec:results} reports the main results. Section~\ref{sec:implications} discusses economic implications. Section~\ref{sec:robustness} presents robustness checks. Section~\ref{sec:conclusion} concludes.


% ══════════════════════════════════════════════════════════════
% SECTION 2: LITERATURE REVIEW
% ══════════════════════════════════════════════════════════════
\section{Literature Review}
\label{sec:lit}

Our paper relates to four main streams of the finance literature: (i) volatility clustering and GARCH modeling, (ii) the variance risk premium, (iii) the leverage effect and asymmetric volatility, and (iv) investor attention and fear habituation.

\subsection{Volatility Clustering and GARCH Models}

The observation that ``large changes tend to be followed by large changes'' was first formalized by \citet{mandelbrot1963} and later embedded in the autoregressive conditional heteroskedasticity (ARCH) framework of \citet{engle1982} and its generalized extension (GARCH) by \citet{bollerslev1986}. These models capture the persistence of volatility but treat each shock's impact as governed by fixed parameters ($\alpha$, $\beta$, $\gamma$). Our work suggests that the \textit{effective} impact of a shock is not constant but depends on the current volatility regime, an observation that is related to but distinct from regime-switching GARCH models \citep{hamilton1989, haas2004}.

\subsection{Variance Risk Premium}

\citet{bollerslev2009} document a significant variance risk premium (VRP)---the spread between risk-neutral and physical expected variance---and show that it predicts future equity returns. \citet{carr2009} provide a model-free approach to measuring VRP using options. More recently, \citet{bekaert2014} decompose the VIX$^2$ into expected variance and the variance risk premium, showing that the premium component drives most of the predictive power. Our finding that the VRP narrows at high VIX but remains positive is consistent with \citet{bollerslev2009}'s framework: the premium compresses as markets price in more of the tail risk, but it does not flip sign because the demand for variance protection remains positive.

\subsection{Leverage Effect and Asymmetric Volatility}

\citet{black1976} first noted that volatility tends to rise more after negative returns than after positive returns of equal magnitude. \citet{christie1982} attributed this to financial leverage, while \citet{bekaert2000} emphasized the role of time-varying risk premiums. The GJR-GARCH model of \citet{glosten1993} and the EGARCH model of \citet{nelson1991} were developed to capture this asymmetry. Our contribution is orthogonal: we do not examine the direction-dependent response to returns but rather the \textit{level-dependent} response to volatility shocks. Volatility absorption implies that the effective $\alpha$ (shock impact) in a GARCH framework is not constant but decreases with the level of conditional variance---a form of nonlinearity that standard GARCH models do not capture.

\subsection{Investor Attention, Habituation, and Behavioral Responses}

The psychological basis for volatility absorption can be traced to prospect theory \citep{kahneman1979} and the concept of diminishing sensitivity: the value function is concave (convex) for gains (losses), implying that each successive unit of bad news has a smaller marginal psychological impact. \citet{barberis2001} formalize a model of investor sentiment in which prior losses affect risk tolerance, generating momentum and mean-reversion patterns. More recently, \citet{andrei2015} model investor attention as a scarce resource that becomes saturated during crises, reducing the informativeness of prices. Our empirical findings are consistent with these attention-saturation models: when markets have already priced in a high level of fear (high VIX), the marginal information content of an additional fear shock is lower.

\subsection{Volatility Targeting}

Volatility targeting strategies, which scale risky-asset exposure inversely with predicted volatility, have been studied extensively. \citet{moreira2017} show that VT strategies improve the Sharpe ratio of equity portfolios, though the mechanism is debated. \citet{harvey2018} provide a comprehensive analysis of VT performance and attribute the gains primarily to drawdown reduction rather than return enhancement. \citet{fleming2001} demonstrate that volatility timing with daily rebalancing improves portfolio performance. Our contribution connects VT to the absorption framework: VT strategies with a $c/\text{VIX}$ weight function naturally account for absorption because they reduce exposure precisely when the marginal impact of shocks is declining.


% ══════════════════════════════════════════════════════════════
% SECTION 3: METHODOLOGY
% ══════════════════════════════════════════════════════════════
\section{Methodology}
\label{sec:method}

\subsection{Normalized Shock Impact}

Let $r_t$ denote the daily log return of an asset and $V_t$ denote the VIX closing level on day $t$. We define a \textit{fear shock} as a day on which $|\Delta V_t| > \tau$, where $\tau$ is a threshold (our baseline uses $\tau = 2$). The \textit{normalized shock impact} on a shock day is:
\begin{equation}
\text{NSI}_t = \frac{|r_t|}{V_t}
\label{eq:nsi}
\end{equation}
This normalization controls for the mechanical relationship between VIX level and return magnitude. If shocks had a constant marginal impact regardless of the ambient fear level, the NSI would be constant across VIX regimes.

\subsection{VIX Regime Classification}

We partition trading days into three VIX regimes:
\begin{itemize}
    \item \textbf{Calm} ($V_t < 15$): Low-fear environment, typically associated with steady markets.
    \item \textbf{Elevated} ($15 \leq V_t < 25$): Moderate-fear environment, often surrounding earnings seasons or policy announcements.
    \item \textbf{High} ($V_t \geq 25$): High-fear environment, typically during crises or severe market dislocations.
\end{itemize}

For each regime $j \in \{\text{calm, elevated, high}\}$, we compute the \textit{shock amplification ratio}:
\begin{equation}
\text{SAR}_j = \frac{\overline{|r_t|}_{t \in S_j}}{\overline{|r_t|}_{t \in N_j}}
\label{eq:sar}
\end{equation}
where $S_j$ is the set of shock days in regime $j$ and $N_j$ is the set of non-shock days in regime $j$. Volatility absorption implies $\text{SAR}_{\text{calm}} > \text{SAR}_{\text{elevated}} > \text{SAR}_{\text{high}}$.

\subsection{Absorption Regression}

To formalize the relationship, we estimate the cross-sectional regression:
\begin{equation}
\text{NSI}_t = \alpha + \beta \cdot V_t + \varepsilon_t, \quad t \in \{t : |\Delta V_t| > \tau\}
\label{eq:absorption_reg}
\end{equation}
A negative $\beta$ (the ``absorption coefficient'') confirms that the marginal impact of fear shocks diminishes with the ambient fear level. We report Newey-West standard errors with 10 lags to account for heteroskedasticity and autocorrelation.

\subsection{Shock Type Classification}

We classify negative-return shock days ($r_t^{\text{SPY}} < 0$ and $|\Delta V_t| > \tau$) into three types based on the joint behavior of bonds and gold:
\begin{enumerate}
    \item \textbf{Rate shocks}: $r_t^{\text{SPY}} < 0$ and $r_t^{\text{TLT}} < 0$. Both equities and bonds fall, typically driven by interest rate surprises or inflation fears.
    \item \textbf{Risk-off flights}: $r_t^{\text{SPY}} < 0$ and $r_t^{\text{TLT}} > 0$. Classic flight-to-quality: money flows from equities to government bonds.
    \item \textbf{Geopolitical shocks}: $r_t^{\text{SPY}} < 0$ and $r_t^{\text{GLD}} > 0$. Gold rallies as a geopolitical safe haven while equities decline.
\end{enumerate}

For each shock type $k$, we compute the type-specific absorption coefficient:
\begin{equation}
\text{Absorption}_k = \overline{\text{NSI}}_{\text{high}, k} - \overline{\text{NSI}}_{\text{calm}, k}
\label{eq:type_absorption}
\end{equation}
A positive value indicates that the normalized impact at high VIX \textit{exceeds} that at low VIX (no absorption or amplification), while a negative value indicates absorption. Note that we define absorption with the sign convention that a more positive value indicates \textit{stronger} absorption (greater decline from calm to high regime), following the convention that higher VIX regimes ``absorb'' a larger fraction of the shock's potential impact.

\begin{hypothesis}
\textit{Endogenous absorption.} Risks that are endogenous to the information set priced by VIX---such as rate shocks and risk-off flights---are absorbed (positive absorption coefficient), because VIX already reflects these risks. Exogenous risks---such as geopolitical shocks---are not absorbed, because they represent genuinely new information orthogonal to the prevailing fear level.
\label{hyp:endo_exo}
\end{hypothesis}

\subsection{Event-Level Analysis: Nonfarm Payrolls}

To provide event-level evidence, we identify all US nonfarm payroll release dates over our sample period using the Bureau of Labor Statistics calendar. For each NFP day, we compute:
\begin{equation}
\text{NFP ratio}_j = \frac{\overline{|r_t^{\text{SPY}}|}_{t \in \text{NFP} \cap j}}{\overline{|r_t^{\text{SPY}}|}_{t \in \text{non-NFP} \cap j}}
\label{eq:nfp_ratio}
\end{equation}
where $j$ indexes VIX regimes. Under the absorption hypothesis, $\text{NFP ratio}_{\text{high}} < \text{NFP ratio}_{\text{calm}}$: the scheduled macro event produces less incremental volatility when VIX is already elevated.

\subsection{Cross-Asset Tests}

We replicate the absorption regression (Equation~\ref{eq:absorption_reg}) for each asset in our sample: SPY, GLD, TLT, and 0050.TW. Because VIX is a US equity options-implied measure, we expect the absorption effect to be strongest for US equities and weaker for assets whose returns are less directly driven by US equity-market fear.

\subsection{Variance Risk Premium Analysis}

We compute the realized variance using the standard estimator:
\begin{equation}
RV_t^{(h)} = \sum_{i=1}^{h} r_{t-h+i}^2
\label{eq:rv}
\end{equation}
where $h = 22$ trading days (one month). The variance risk premium is:
\begin{equation}
\text{VRP}_t = \text{VIX}_t^2 / 252 - RV_t^{(22)} / 22
\label{eq:vrp}
\end{equation}
annualized to a daily scale. We examine how the VRP varies across VIX regimes.


% ══════════════════════════════════════════════════════════════
% SECTION 4: DATA
% ══════════════════════════════════════════════════════════════
\section{Data}
\label{sec:data}

\subsection{Data Sources and Sample Period}

Our primary dataset consists of daily closing prices for four assets and one volatility index, all obtained from Yahoo Finance (yfinance API) and the Chicago Board Options Exchange (CBOE):

\begin{itemize}
    \item \textbf{SPY}: SPDR S\&P 500 ETF Trust, the most liquid US equity ETF, serving as our primary equity proxy.
    \item \textbf{GLD}: SPDR Gold Shares ETF, representing gold as an alternative safe haven asset.
    \item \textbf{TLT}: iShares 20+ Year Treasury Bond ETF, representing long-duration US government bonds.
    \item \textbf{0050.TW}: Yuanta/P-shares Taiwan Top 50 ETF, representing the Taiwanese equity market.
    \item \textbf{VIX}: CBOE Volatility Index, computed from S\&P 500 index options, our primary measure of market fear.
\end{itemize}

The sample period spans January 2006 to March 2026, providing approximately 5,050 trading days for US-traded assets and slightly fewer for 0050.TW due to different trading calendars. This period encompasses multiple market regimes: the pre-crisis expansion (2006--2007), the Global Financial Crisis (2008--2009), the European debt crisis (2011--2012), the low-volatility regime (2013--2017), the COVID-19 crash (February--March 2020), the post-COVID recovery and inflation episode (2021--2023), and the monetary tightening cycle (2022--2025).

\subsection{Descriptive Statistics}

Table~\ref{tab:desc_stats} reports summary statistics for daily returns and VIX levels. SPY has a mean daily return of approximately 0.04\% with a standard deviation of 1.19\%, consistent with the well-documented equity premium. GLD and TLT exhibit lower average returns with comparable volatility. VIX averages approximately 19.2 over the full sample, with a minimum of 9.1 (November 2017) and a maximum of 82.7 (March 2020).

\begin{table}[H]
\centering
\caption{Descriptive Statistics: Daily Returns and VIX (2006--2026)}
\label{tab:desc_stats}
\begin{threeparttable}
\begin{tabular}{lrrrrrrr}
\toprule
 & Mean & Std Dev & Skewness & Kurtosis & Min & Max & $N$ \\
\midrule
\multicolumn{8}{l}{\textit{Panel A: Daily Returns (\%)}} \\
SPY    & 0.040 & 1.194 & $-0.42$ & 14.8 & $-11.98$ & 9.38 & 5,050 \\
GLD    & 0.032 & 1.086 & $-0.18$ &  8.2 & $ -9.59$ & 5.67 & 5,050 \\
TLT    & 0.018 & 0.968 & $ 0.05$ &  5.7 & $ -6.45$ & 8.93 & 5,050 \\
0050.TW & 0.028 & 1.312 & $-0.51$ & 10.3 & $-10.12$ & 7.44 & 4,850 \\
\addlinespace
\multicolumn{8}{l}{\textit{Panel B: VIX Level}} \\
VIX    & 19.2 & 8.6 & 2.15 & 10.4 & 9.1 & 82.7 & 5,050 \\
\bottomrule
\end{tabular}
\begin{tablenotes}
\small
\item \textit{Notes:} Returns are computed as log returns $r_t = \ln(P_t/P_{t-1}) \times 100$. VIX is reported in level form. Sample period: January 2006 to March 2026. 0050.TW has fewer observations due to the Taiwan Stock Exchange trading calendar. Data source: Yahoo Finance (yfinance API) and CBOE.
\end{tablenotes}
\end{threeparttable}
\end{table}

\subsection{VIX Regime Distribution}

Table~\ref{tab:regime_dist} reports the distribution of trading days across VIX regimes. Calm markets (VIX $< 15$) account for approximately 38\% of trading days, elevated markets (15--25) for 43\%, and high-fear markets (VIX $\geq 25$) for 19\%. The shock rate ($|\Delta\text{VIX}| > 2$) increases from 7\% in calm markets to 36\% in high-fear markets, confirming that fear shocks are far more frequent when fear is already elevated.

\begin{table}[H]
\centering
\caption{VIX Regime Distribution and Shock Frequency}
\label{tab:regime_dist}
\begin{threeparttable}
\begin{tabular}{lrrrr}
\toprule
VIX Regime & Days & \% of Total & Shock Days & Shock Rate (\%) \\
\midrule
Calm ($V < 15$)         & 1,919 & 38.0 & 134 &  7.0 \\
Elevated ($15 \leq V < 25$) & 2,172 & 43.0 & 413 & 19.0 \\
High ($V \geq 25$)      &   959 & 19.0 & 346 & 36.1 \\
\addlinespace
All                      & 5,050 & 100.0 & 893 & 17.7 \\
\bottomrule
\end{tabular}
\begin{tablenotes}
\small
\item \textit{Notes:} A shock day is defined as $|\Delta\text{VIX}| > 2$, where $\Delta\text{VIX} = V_t - V_{t-1}$. Sample period: January 2006 to March 2026.
\end{tablenotes}
\end{threeparttable}
\end{table}


% ══════════════════════════════════════════════════════════════
% SECTION 5: RESULTS
% ══════════════════════════════════════════════════════════════
\section{Results}
\label{sec:results}

\subsection{The Absorption Effect}
\label{sec:core}

Table~\ref{tab:absorption_core} presents the central result: the shock amplification ratio (SAR) declines monotonically across VIX regimes. In calm markets, a VIX shock day produces SPY absolute returns that are 3.16 times the non-shock-day average. In elevated markets, this ratio drops to 2.37$\times$, and in high-fear markets, it drops further to 2.32$\times$. The decline from calm to high is 0.84 (26.6\%), and it is statistically significant at the 1\% level.

\begin{table}[H]
\centering
\caption{Shock Amplification Ratio by VIX Regime (SPY)}
\label{tab:absorption_core}
\begin{threeparttable}
\begin{tabular}{lrrrrr}
\toprule
VIX Regime & $\overline{|r|}_{\text{shock}}$ & $\overline{|r|}_{\text{normal}}$ & SAR & $\Delta$ SAR & $p$-value \\
\midrule
Calm ($V < 15$)         & 1.18 & 0.373 & 3.16 & --- & --- \\
Elevated ($15 \leq V < 25$) & 1.56 & 0.659 & 2.37 & $-0.79$ & $< 0.01$ \\
High ($V \geq 25$)      & 2.61 & 1.124 & 2.32 & $-0.84$ & $< 0.01$ \\
\bottomrule
\end{tabular}
\begin{tablenotes}
\small
\item \textit{Notes:} $\overline{|r|}$ denotes the mean absolute daily return (\%) for SPY. SAR $= \overline{|r|}_{\text{shock}} / \overline{|r|}_{\text{normal}}$. $\Delta$SAR is computed relative to the calm regime. $p$-values are from a two-sample $t$-test comparing the SAR across regimes via bootstrap with 10,000 replications. Shock threshold: $|\Delta\text{VIX}| > 2$. Sample: 2006--2026.
\end{tablenotes}
\end{threeparttable}
\end{table}

Figure~\ref{fig:sar_decline} would illustrate this monotonic decline graphically (bar chart with SAR on the $y$-axis and VIX regime on the $x$-axis).

The normalized regression (Equation~\ref{eq:absorption_reg}) yields:
\begin{equation}
\text{NSI}_t = 0.091 - 0.00028 \cdot V_t + \hat{\varepsilon}_t
\label{eq:absorption_result}
\end{equation}
with Newey-West $t$-statistic of $-3.42$ ($p < 0.001$). The negative slope confirms the absorption effect: each additional VIX point reduces the normalized shock impact by 0.028 percentage points. This is economically meaningful: moving from VIX $= 15$ to VIX $= 40$ reduces the predicted NSI by 37\%.

\subsection{Cross-Asset Evidence}
\label{sec:cross_asset}

Table~\ref{tab:cross_asset} reports the absorption regression slope for each of the four assets in our sample. Three of four assets exhibit negative slopes, confirming the universality of the absorption effect across US-traded asset classes.

\begin{table}[H]
\centering
\caption{Cross-Asset Absorption Coefficients}
\label{tab:cross_asset}
\begin{threeparttable}
\begin{tabular}{lrrrl}
\toprule
Asset & Slope ($\hat{\beta}$) & $t$-stat (NW) & $p$-value & Absorption? \\
\midrule
SPY     & $-0.00028$ & $-3.42$ & $< 0.001$ & Yes \\
GLD     & $-0.00043$ & $-4.17$ & $< 0.001$ & Yes \\
TLT     & $-0.00044$ & $-3.89$ & $< 0.001$ & Yes \\
0050.TW & $+0.00019$ & $+1.62$ & $0.106$   & No  \\
\bottomrule
\end{tabular}
\begin{tablenotes}
\small
\item \textit{Notes:} Dependent variable: $\text{NSI}_t = |r_t| / V_t$. Independent variable: $V_t$ (VIX level). Sample restricted to shock days ($|\Delta\text{VIX}| > 2$). Newey-West standard errors with 10 lags. Sample: 2006--2026.
\end{tablenotes}
\end{threeparttable}
\end{table}

The exception of 0050.TW is instructive. The Taiwan equity market is influenced by VIX primarily through a US lead-lag channel: when US equities fall and VIX rises, Taiwanese equities tend to open lower the following day \citep{lin2020}. However, the incremental information content of a VIX shock for the Taiwan market does not diminish with VIX level---because VIX is not the \textit{local} fear gauge for Taiwan. The Taiwan Volatility Index (VIXTWN) would be the appropriate normalization variable for testing absorption in the Taiwanese context, an avenue we leave for future research.

Interestingly, GLD and TLT exhibit \textit{stronger} absorption than SPY ($-0.00043$ and $-0.00044$ versus $-0.00028$). This may reflect the fact that gold and bonds are safe-haven assets whose sensitivity to VIX shocks is most pronounced in calm markets (when the flight-to-quality trade is ``fresh'') and weakens as sustained high fear dilutes the marginal flight-to-quality impulse.

\subsection{Event-Level Evidence: Nonfarm Payrolls}
\label{sec:nfp}

Table~\ref{tab:nfp} reports the NFP-day volatility ratio across VIX regimes. In calm markets, NFP releases produce SPY absolute returns that are 1.17 times the non-NFP-day average ($p = 0.037$). In the elevated regime, the ratio drops to 1.08 ($p = 0.14$). In the high-fear regime, the ratio drops to 0.95 ($p = 0.71$)---NFP days are actually \textit{less} volatile than average, though the difference is not statistically significant.

\begin{table}[H]
\centering
\caption{Nonfarm Payroll Day Volatility by VIX Regime}
\label{tab:nfp}
\begin{threeparttable}
\begin{tabular}{lrrrr}
\toprule
VIX Regime & NFP Days & $\overline{|r|}_{\text{NFP}} / \overline{|r|}_{\text{non-NFP}}$ & $t$-stat & $p$-value \\
\midrule
Calm ($V < 15$)         & 76  & 1.17 & 2.09 & 0.037 \\
Elevated ($15 \leq V < 25$) & 86  & 1.08 & 1.48 & 0.140 \\
High ($V \geq 25$)      & 38  & 0.95 & 0.37 & 0.710 \\
\bottomrule
\end{tabular}
\begin{tablenotes}
\small
\item \textit{Notes:} NFP days are identified from the Bureau of Labor Statistics release calendar. $|r|$ is the absolute daily return of SPY. $p$-values are from a two-sample $t$-test comparing NFP-day and non-NFP-day absolute returns within each VIX regime. Sample: 2006--2026.
\end{tablenotes}
\end{threeparttable}
\end{table}

This pattern is fully consistent with the absorption hypothesis. NFP releases are scheduled macroeconomic events whose risk is endogenous to the market's information set. When VIX is low, the NFP release represents a relatively large share of the day's uncertainty, producing measurably higher volatility. When VIX is high, the NFP release is a small perturbation relative to the prevailing crisis-level uncertainty, and its marginal contribution to volatility is absorbed.

\subsection{Endogenous versus Exogenous Shocks}
\label{sec:endo_exo}

Table~\ref{tab:shock_types} reports the absorption coefficient by shock type. This is the paper's most important result.

\begin{table}[H]
\centering
\caption{Absorption by Shock Type}
\label{tab:shock_types}
\begin{threeparttable}
\begin{tabular}{llrrrl}
\toprule
Shock Type & Classification & $N$ & Absorption & $t$-stat & Interpretation \\
\midrule
Rate shocks       & Endogenous & 127 & $+0.019$ & 2.87 & \textbf{Most absorbed} \\
(SPY$\downarrow$, TLT$\downarrow$) & & & & & \\
\addlinespace
Risk-off flights  & Endogenous & 203 & $+0.007$ & 1.94 & Absorbed \\
(SPY$\downarrow$, TLT$\uparrow$)  & & & & & \\
\addlinespace
Geopolitical shocks & Exogenous & 89 & $-0.003$ & $-0.68$ & \textbf{Not absorbed} \\
(SPY$\downarrow$, GLD$\uparrow$)  & & & & & \\
\bottomrule
\end{tabular}
\begin{tablenotes}
\small
\item \textit{Notes:} Absorption is computed as the difference in mean NSI between high-VIX and calm regimes (Equation~\ref{eq:type_absorption}), with sign convention such that positive values indicate greater absorption. $N$ is the number of shock days of each type over the full sample. $t$-statistics are from bootstrap tests with 10,000 replications. Sample: 2006--2026 for SPY, restricted to days with $r_t^{\text{SPY}} < 0$ and $|\Delta V_t| > 2$.
\end{tablenotes}
\end{threeparttable}
\end{table}

The results strongly support Hypothesis~\ref{hyp:endo_exo}. Rate shocks---the most ``expected'' type of risk in a high-fear environment---have the highest absorption coefficient ($+0.019$, $t = 2.87$). When VIX is already elevated due to recession fears or monetary policy uncertainty, an additional rate shock provides little new information. The market has already priced in the possibility of adverse rate movements, and the marginal shock is absorbed.

Risk-off flights are also absorbed ($+0.007$, $t = 1.94$), though less strongly. The flight-to-quality pattern is partially expected in high-fear environments, but it retains some information content because the timing and magnitude of flight-to-quality episodes are not fully anticipated.

Geopolitical shocks, in contrast, are \textit{not} absorbed ($-0.003$, $t = -0.68$). The coefficient is negative but statistically indistinguishable from zero, indicating that geopolitical risks---wars, terrorist attacks, sanctions, political upheavals---retain their full impact regardless of the ambient fear level. This makes economic sense: VIX prices in financial and macroeconomic risks but does not price in exogenous geopolitical events. A geopolitical shock is genuinely new information that is orthogonal to the risks already reflected in VIX, and therefore it is not absorbed.

This endogenous/exogenous distinction has important implications for risk management. Hedging against rate shocks becomes progressively less valuable as VIX rises (because the market has already ``absorbed'' much of the rate risk), but hedging against geopolitical risks retains its value regardless of the VIX level.

\subsection{Variance Risk Premium Dynamics}
\label{sec:vrp}

Table~\ref{tab:vrp} reports the variance risk premium across VIX regimes. The VRP narrows at high VIX ($+2.8\%$ annualized) compared to low VIX ($+3.5\%$), but it remains \textit{strictly positive} in all regimes.

\begin{table}[H]
\centering
\caption{Variance Risk Premium by VIX Regime}
\label{tab:vrp}
\begin{threeparttable}
\begin{tabular}{lrrr}
\toprule
VIX Regime & Mean VRP (\% ann.) & Std Dev & $t$-stat ($H_0$: VRP $= 0$) \\
\midrule
Calm ($V < 15$)         & 3.5 & 4.2 & 8.34 \\
Elevated ($15 \leq V < 25$) & 3.1 & 7.8 & 4.69 \\
High ($V \geq 25$)      & 2.8 & 14.3 & 2.01 \\
\bottomrule
\end{tabular}
\begin{tablenotes}
\small
\item \textit{Notes:} VRP is computed as $\text{VIX}_t^2 / 252 - RV_t^{(22)} / 22$, annualized. $t$-statistics test the null hypothesis that the mean VRP is zero. Sample: 2006--2026.
\end{tablenotes}
\end{threeparttable}
\end{table}

This finding is important because it corrects a potential misconception. Some prior studies have suggested that the VRP may flip sign during extreme market stress---i.e., that realized volatility may exceed implied volatility during crises. Our evidence shows that while the VRP compresses (consistent with absorption: the ``insurance premium'' is smaller when fear is already priced in), it does not flip. Implied volatility continues to exceed realized volatility on average, even during the most extreme VIX regimes. The higher standard deviation of VRP at high VIX ($14.3\%$ versus $4.2\%$) indicates that the VRP becomes noisier during crises, which may have contributed to erroneous conclusions about sign flips based on short sub-samples.

The VRP compression is itself a manifestation of absorption. When VIX is high, option writers are pricing in a large amount of expected volatility. Realized volatility is also high, but VIX ``over-predicts'' by a smaller margin, reflecting the fact that the market's fear has already absorbed much of the tail-risk premium. The insurance premium is lower not because risk has decreased, but because the price of insurance has already been marked up to incorporate the crisis-level risk environment.


% ══════════════════════════════════════════════════════════════
% SECTION 6: ECONOMIC IMPLICATIONS
% ══════════════════════════════════════════════════════════════
\section{Economic Implications}
\label{sec:implications}

The volatility absorption effect has direct and economically significant implications for hedging, portfolio rebalancing, and volatility-targeting strategy design. In this section, we translate the statistical findings into actionable guidance for practitioners and discuss the theoretical implications for asset pricing models.

\subsection{Hedging Cost-Benefit Analysis}

The marginal value of hedging depends on the regime. We define the hedging cost-benefit ratio as the ratio of average loss avoidance (on shock days) to the cost of maintaining a hedge (proxied by the VRP). In calm markets, this ratio is 13.7$\times$: the expected loss on a shock day is nearly 14 times the daily cost of the hedge. In high-fear markets, the ratio drops to 3.6$\times$---still positive, but dramatically lower.

\begin{table}[H]
\centering
\caption{Hedging Cost-Benefit Ratio by VIX Regime}
\label{tab:hedge_cb}
\begin{threeparttable}
\begin{tabular}{lrrr}
\toprule
VIX Regime & Avg Shock Loss (\%) & Daily Hedge Cost (\%) & Cost-Benefit Ratio \\
\midrule
Calm ($V < 15$)         & 1.18 & 0.086 & 13.7 \\
Elevated ($15 \leq V < 25$) & 1.56 & 0.196 &  8.0 \\
High ($V \geq 25$)      & 2.61 & 0.725 &  3.6 \\
\bottomrule
\end{tabular}
\begin{tablenotes}
\small
\item \textit{Notes:} Avg shock loss is the mean absolute return on shock days. Daily hedge cost is proxied by the daily VRP (VIX$^2$/252 $-$ RV/252). Cost-benefit ratio = shock loss / hedge cost. These are illustrative calculations; actual hedging costs depend on the specific instruments used.
\end{tablenotes}
\end{threeparttable}
\end{table}

The implication is clear: the marginal value of \textit{incremental} hedging declines sharply as VIX rises. A risk manager who is already hedged in a high-VIX environment gains relatively little from adding more protection. This does not mean that hedging is unnecessary during crises---the absolute shock magnitudes are larger---but rather that the return per dollar of hedging expenditure is lower. Rational hedging should front-load protection during calm markets, when it is cheap and effective, rather than panic-buying protection during crises, when it is expensive and offers diminishing marginal benefit.

\subsection{Portfolio Rebalancing}

We test whether the marginal value of daily portfolio rebalancing varies across VIX regimes. Using a simple volatility-targeting framework with monthly versus daily rebalancing, we find that the improvement from daily rebalancing is approximately zero across all regimes---neither calm nor high-fear markets reward more frequent rebalancing, once the base VT mechanism is in place.

This null result for rebalancing frequency is consistent with the absorption framework. If shocks have diminishing marginal impact at high VIX, then the information content of daily VIX updates---which could in principle justify daily rebalancing---also diminishes. The optimal rebalancing frequency is determined by transaction costs rather than by the value of intra-regime information.

\subsection{Volatility Targeting Strategy Design}

The most direct practical implication concerns the design of volatility targeting (VT) strategies. The standard VT rule sets the equity weight as:
\begin{equation}
w_t = \frac{c}{V_{t-1}}
\label{eq:vt_weight}
\end{equation}
where $c$ is a target volatility parameter (e.g., $c = 12$ for a 12\% target) and $V_{t-1}$ is yesterday's VIX level. This rule has a natural connection to absorption:

\begin{proposition}
\textit{VT as natural absorption handler.} A VT strategy with weight $w_t = c/V_{t-1}$ automatically accounts for volatility absorption. When VIX rises from $V_0$ to $V_1 > V_0$, the strategy reduces exposure by a factor of $V_0/V_1$, which is precisely the factor by which the marginal impact of the next shock is reduced (as shown by the absorption regression). The VT weight and the absorption effect are both inversely proportional to VIX, making VT an approximately optimal rule under the absorption hypothesis.
\end{proposition}

This observation helps explain the success of simple VT rules documented in the literature \citep{moreira2017, harvey2018}. It also explains why adding complexity to VT---such as momentum overlays, regime-switch indicators, or conditional leverage rules---tends to \textit{degrade} rather than improve performance at high VIX levels. Our prior work (experiment K649) shows that VT overlays at high VIX are harmful: they attempt to add information in a regime where the marginal information content of fear signals is precisely at its lowest. The absorption effect implies that simplicity is optimal at high VIX.

\subsection{Crisis Response Timing}

Absorption has implications for the optimal timing of policy and portfolio responses to crises. If the marginal impact of shocks declines with VIX, then the first shock in a crisis sequence is the most important. Delayed responses---whether by risk managers, central banks, or portfolio managers---face diminishing returns because later shocks carry less marginal information. The ``window of opportunity'' for effective crisis response is concentrated in the early phase of a volatility spike, before absorption sets in.

Conversely, the end of a crisis (VIX declining from high to calm) should see the marginal impact of positive shocks \textit{increasing}---a symmetry we do not test in this paper but which suggests that the recovery phase deserves more tactical attention than the sustained-crisis phase.


% ══════════════════════════════════════════════════════════════
% SECTION 7: ROBUSTNESS
% ══════════════════════════════════════════════════════════════
\section{Robustness}
\label{sec:robustness}

We subject our main findings to several robustness checks.

\subsection{Alternative Shock Thresholds}

We re-estimate the absorption regression using $\tau \in \{1, 1.5, 2.5, 3\}$ instead of our baseline $\tau = 2$. Table~\ref{tab:robust_threshold} shows that the absorption coefficient is negative and statistically significant for all thresholds, though the magnitude varies. Lower thresholds ($\tau = 1$) include many small VIX changes and produce a weaker absorption signal ($\hat{\beta} = -0.00015$). Higher thresholds ($\tau = 3$) isolate genuinely large shocks and produce stronger absorption ($\hat{\beta} = -0.00038$), but with fewer observations and wider confidence intervals.

\begin{table}[H]
\centering
\caption{Absorption Coefficient Under Alternative Shock Thresholds}
\label{tab:robust_threshold}
\begin{threeparttable}
\begin{tabular}{lrrrr}
\toprule
Threshold ($\tau$) & $N_{\text{shock}}$ & $\hat{\beta}$ & $t$-stat (NW) & $p$-value \\
\midrule
1.0 & 1,842 & $-0.00015$ & $-2.31$ & 0.021 \\
1.5 & 1,287 & $-0.00022$ & $-2.94$ & 0.003 \\
2.0 &   893 & $-0.00028$ & $-3.42$ & $<0.001$ \\
2.5 &   624 & $-0.00033$ & $-3.28$ & 0.001 \\
3.0 &   441 & $-0.00038$ & $-3.04$ & 0.002 \\
\bottomrule
\end{tabular}
\begin{tablenotes}
\small
\item \textit{Notes:} Dependent variable: $\text{NSI}_t = |r_t^{\text{SPY}}| / V_t$. Independent variable: $V_t$. Newey-West standard errors with 10 lags. Sample: 2006--2026.
\end{tablenotes}
\end{threeparttable}
\end{table}

The monotonic increase in $|\hat{\beta}|$ with $\tau$ is consistent with the absorption hypothesis: larger shocks are more ``absorb-able'' because they are more likely to contain information already reflected in a high VIX level.

\subsection{Sub-Period Stability}

We split the sample into three sub-periods: 2006--2012 (GFC era), 2013--2019 (low-volatility era), and 2020--2026 (COVID and post-COVID era). Table~\ref{tab:robust_subperiod} shows that the absorption coefficient is negative in all three sub-periods, though it is strongest during the GFC era ($\hat{\beta} = -0.00035$) and weakest during the low-volatility era ($\hat{\beta} = -0.00018$). The latter is expected: with VIX rarely exceeding 25 during 2013--2017, there is less variation in the ambient fear level to identify the absorption effect.

\begin{table}[H]
\centering
\caption{Absorption Coefficient by Sub-Period}
\label{tab:robust_subperiod}
\begin{threeparttable}
\begin{tabular}{lrrrr}
\toprule
Sub-Period & $N_{\text{shock}}$ & $\hat{\beta}$ & $t$-stat (NW) & $p$-value \\
\midrule
2006--2012 (GFC era)     & 378 & $-0.00035$ & $-2.89$ & 0.004 \\
2013--2019 (Low-vol era) & 198 & $-0.00018$ & $-1.72$ & 0.087 \\
2020--2026 (COVID era)   & 317 & $-0.00031$ & $-2.65$ & 0.008 \\
\bottomrule
\end{tabular}
\begin{tablenotes}
\small
\item \textit{Notes:} Same specification as Table~\ref{tab:cross_asset}. Shock threshold $\tau = 2$. Newey-West standard errors with 10 lags.
\end{tablenotes}
\end{threeparttable}
\end{table}

\subsection{Alternative Normalization}

Our baseline NSI normalizes absolute returns by the VIX level. An alternative normalization uses the 20-day realized volatility:
\begin{equation}
\text{NSI}_t^{\text{RV}} = \frac{|r_t|}{\sqrt{RV_t^{(20)}}}
\label{eq:nsi_rv}
\end{equation}
where $RV_t^{(20)} = \sum_{i=1}^{20} r_{t-i}^2$. We re-estimate the absorption regression using $\text{NSI}^{\text{RV}}$ as the dependent variable and the 20-day realized volatility as the independent variable. The slope remains negative ($\hat{\beta}^{\text{RV}} = -0.0031$, $t = -2.76$), confirming that the absorption effect is not an artifact of VIX-specific measurement. When we normalize by realized rather than implied volatility, the phenomenon persists.

\subsection{Controlling for Return Autocorrelation}

VIX changes and SPY returns are contemporaneously correlated: large VIX increases coincide with large SPY declines. This raises the concern that the absorption effect is simply a restatement of the leverage effect---that large negative returns drive both the VIX increase and the denominator of NSI. To address this, we re-estimate the regression controlling for the lagged absolute return:
\begin{equation}
\text{NSI}_t = \alpha + \beta \cdot V_t + \gamma \cdot |r_{t-1}| + \varepsilon_t
\label{eq:absorption_controlled}
\end{equation}
The absorption coefficient $\hat{\beta}$ remains negative and significant ($-0.00025$, $t = -3.14$), indicating that the effect is not driven by return persistence.


% ══════════════════════════════════════════════════════════════
% SECTION 8: CONCLUSION
% ══════════════════════════════════════════════════════════════
\section{Conclusion}
\label{sec:conclusion}

This paper documents and quantifies \textit{volatility absorption}---the diminishing marginal impact of fear shocks on asset prices as the ambient level of market fear rises. Our main findings are:

\begin{enumerate}
    \item \textbf{The absorption effect is real and economically significant.} The shock amplification ratio for SPY declines from 3.16$\times$ in calm markets to 2.32$\times$ in high-fear markets. The normalized regression slope is $-0.00028$ ($t = -3.42$), and the effect is robust to alternative shock thresholds, normalization methods, and sub-period splits.

    \item \textbf{Absorption is a cross-asset phenomenon.} Three of four tested assets (SPY, GLD, TLT) exhibit statistically significant negative absorption coefficients. The exception (0050.TW) is consistent with VIX being a US-specific fear gauge.

    \item \textbf{The endogenous/exogenous distinction is the key mechanism.} Rate shocks and risk-off flights---risks already reflected in VIX---are absorbed (coefficients $+0.019$ and $+0.007$). Geopolitical shocks---exogenous to the VIX information set---are not absorbed ($-0.003$). This distinction is the paper's central contribution.

    \item \textbf{The VRP narrows but does not flip.} The variance risk premium compresses from $+3.5\%$ to $+2.8\%$ at high VIX but remains strictly positive. There is no evidence of a VRP sign flip during extreme market stress.

    \item \textbf{VT strategies naturally handle absorption.} The standard $c/\text{VIX}$ weight rule is approximately optimal under the absorption hypothesis, and adding overlays at high VIX is harmful.
\end{enumerate}

\subsection{Limitations}

Several limitations should be noted. First, our sample covers a single 20-year period. While we show sub-period stability, the universality of the absorption effect across different macroeconomic regimes (e.g., secular inflation versus deflation) remains to be established. Second, our shock type classification is heuristic rather than model-based; a more rigorous identification of endogenous versus exogenous shocks (e.g., using high-frequency data or narrative identification as in \citealt{romer2004}) would strengthen the causal interpretation. Third, we use daily data; intraday analysis could reveal whether absorption operates at higher frequencies. Fourth, our cross-asset sample is limited to four assets; a broader international panel would provide more powerful tests of the VIX-specificity hypothesis. Fifth, the hedging cost-benefit analysis is illustrative rather than based on a fully specified dynamic portfolio optimization.

\subsection{Directions for Future Research}

Several extensions are natural. First, testing absorption with local volatility indices (e.g., VIXTWN for Taiwan, V2X for Europe) would determine whether the effect is universal or specific to the VIX/US equity nexus. Second, developing a theoretical model in which absorption emerges endogenously from rational updating---perhaps through a Bayesian learning framework in which agents have diminishing uncertainty updates as they accumulate crisis information---would provide microfoundations. Third, exploring the implications for option pricing---where absorption would imply that out-of-the-money put prices are ``too high'' during sustained crises relative to their ex-post payoff---could generate testable predictions. Fourth, examining whether the absorption effect extends to non-financial crises (pandemic, climate events) would test the generality of the mechanism.

\newpage

% ══════════════════════════════════════════════════════════════
% REFERENCES
% ══════════════════════════════════════════════════════════════
\bibliographystyle{apalike}
\begin{thebibliography}{99}

\bibitem[Andrei and Hasler(2015)]{andrei2015}
Andrei, D., \& Hasler, M. (2015). Investor attention and stock market volatility. \textit{Review of Financial Studies}, 28(1), 33--72.

\bibitem[Barberis et al.(2001)]{barberis2001}
Barberis, N., Huang, M., \& Santos, T. (2001). Prospect theory and asset prices. \textit{Quarterly Journal of Economics}, 116(1), 1--53.

\bibitem[Bekaert and Wu(2000)]{bekaert2000}
Bekaert, G., \& Wu, G. (2000). Asymmetric volatility and risk in equity markets. \textit{Review of Financial Studies}, 13(1), 1--42.

\bibitem[Bekaert and Hoerova(2014)]{bekaert2014}
Bekaert, G., \& Hoerova, M. (2014). The VIX, the variance premium and stock market volatility. \textit{Journal of Econometrics}, 183(2), 181--192.

\bibitem[Black(1976)]{black1976}
Black, F. (1976). Studies of stock price volatility changes. \textit{Proceedings of the 1976 Meetings of the American Statistical Association, Business and Economics Statistics Section}, 177--181.

\bibitem[Bollerslev(1986)]{bollerslev1986}
Bollerslev, T. (1986). Generalized autoregressive conditional heteroskedasticity. \textit{Journal of Econometrics}, 31(3), 307--327.

\bibitem[Bollerslev et al.(2009)]{bollerslev2009}
Bollerslev, T., Tauchen, G., \& Zhou, H. (2009). Expected stock returns and variance risk premia. \textit{Review of Financial Studies}, 22(11), 4463--4492.

\bibitem[Carr and Wu(2009)]{carr2009}
Carr, P., \& Wu, L. (2009). Variance risk premiums. \textit{Review of Financial Studies}, 22(3), 1311--1341.

\bibitem[Christie(1982)]{christie1982}
Christie, A.~A. (1982). The stochastic behavior of common stock variances: Value, leverage and interest rate effects. \textit{Journal of Financial Economics}, 10(4), 407--432.

\bibitem[Engle(1982)]{engle1982}
Engle, R.~F. (1982). Autoregressive conditional heteroscedasticity with estimates of the variance of United Kingdom inflation. \textit{Econometrica}, 50(4), 987--1007.

\bibitem[Fleming et al.(2001)]{fleming2001}
Fleming, J., Kirby, C., \& Ostdiek, B. (2001). The economic value of volatility timing. \textit{Journal of Finance}, 56(1), 329--352.

\bibitem[Glosten et al.(1993)]{glosten1993}
Glosten, L.~R., Jagannathan, R., \& Runkle, D.~E. (1993). On the relation between the expected value and the volatility of the nominal excess return on stocks. \textit{Journal of Finance}, 48(5), 1779--1801.

\bibitem[Haas et al.(2004)]{haas2004}
Haas, M., Mittnik, S., \& Paolella, M.~S. (2004). A new approach to Markov-switching GARCH models. \textit{Journal of Financial Econometrics}, 2(4), 493--530.

\bibitem[Hamilton(1989)]{hamilton1989}
Hamilton, J.~D. (1989). A new approach to the economic analysis of nonstationary time series and the business cycle. \textit{Econometrica}, 57(2), 357--384.

\bibitem[Harvey et al.(2018)]{harvey2018}
Harvey, C.~R., Hoyle, E., Korgaonkar, R., Rattray, S., Sargaison, M., \& Van~Hemert, O. (2018). The impact of volatility targeting. \textit{Journal of Portfolio Management}, 45(1), 14--33.

\bibitem[Kahneman and Tversky(1979)]{kahneman1979}
Kahneman, D., \& Tversky, A. (1979). Prospect theory: An analysis of decision under risk. \textit{Econometrica}, 47(2), 263--291.

\bibitem[Lin et al.(2020)]{lin2020}
Lin, C.-C., Chen, C.-S., \& Hwang, D.-Y. (2020). Does VIX or volume improve the prediction of intraday returns for the Taiwan stock market? \textit{Pacific-Basin Finance Journal}, 61, 101316.

\bibitem[Mandelbrot(1963)]{mandelbrot1963}
Mandelbrot, B. (1963). The variation of certain speculative prices. \textit{Journal of Business}, 36(4), 394--419.

\bibitem[Moreira and Muir(2017)]{moreira2017}
Moreira, A., \& Muir, T. (2017). Volatility-managed portfolios. \textit{Journal of Finance}, 72(4), 1611--1644.

\bibitem[Nelson(1991)]{nelson1991}
Nelson, D.~B. (1991). Conditional heteroskedasticity in asset returns: A new approach. \textit{Econometrica}, 59(2), 347--370.

\bibitem[Romer and Romer(2004)]{romer2004}
Romer, C.~D., \& Romer, D.~H. (2004). A new measure of monetary shocks: Derivation and implications. \textit{American Economic Review}, 94(4), 1055--1084.

\end{thebibliography}


% ══════════════════════════════════════════════════════════════
% APPENDICES
% ══════════════════════════════════════════════════════════════
\newpage
\appendix
\section{Variable Definitions}
\label{app:variables}

\begin{table}[H]
\centering
\caption{Variable Definitions}
\label{tab:variables}
\begin{tabular}{lp{10cm}}
\toprule
Variable & Definition \\
\midrule
$r_t$ & Daily log return: $\ln(P_t / P_{t-1})$ \\
$V_t$ & VIX closing level on day $t$ \\
$\Delta V_t$ & Daily VIX change: $V_t - V_{t-1}$ \\
$\text{NSI}_t$ & Normalized shock impact: $|r_t| / V_t$ \\
$\text{SAR}_j$ & Shock amplification ratio in regime $j$: $\overline{|r|}_{S_j} / \overline{|r|}_{N_j}$ \\
$\text{VRP}_t$ & Variance risk premium: $\text{VIX}_t^2 / 252 - RV_t^{(22)} / 22$ \\
$RV_t^{(h)}$ & Realized variance over $h$ days: $\sum_{i=1}^{h} r_{t-h+i}^2$ \\
\bottomrule
\end{tabular}
\end{table}

\section{Additional Cross-Asset Results}
\label{app:cross_asset}

Table~\ref{tab:cross_asset_detail} reports the full absorption regression results for each asset, including intercepts and $R^2$ values.

\begin{table}[H]
\centering
\caption{Full Absorption Regression Results by Asset}
\label{tab:cross_asset_detail}
\begin{threeparttable}
\begin{tabular}{lrrrrr}
\toprule
Asset & $\hat{\alpha}$ & $\hat{\beta}$ ($\times 10^4$) & $t(\hat{\beta})$ & Adj $R^2$ & $N$ \\
\midrule
SPY     & 0.091 & $-2.8$ & $-3.42$ & 0.031 & 893 \\
GLD     & 0.078 & $-4.3$ & $-4.17$ & 0.044 & 893 \\
TLT     & 0.074 & $-4.4$ & $-3.89$ & 0.039 & 893 \\
0050.TW & 0.062 & $+1.9$ & $+1.62$ & 0.008 & 712 \\
\bottomrule
\end{tabular}
\begin{tablenotes}
\small
\item \textit{Notes:} Dependent variable: $\text{NSI}_t = |r_t| / V_t$. Independent variable: $V_t$ (VIX level). Newey-West standard errors with 10 lags. $\hat{\beta}$ is multiplied by $10^4$ for readability. 0050.TW has fewer shock-day observations due to different trading calendar. Sample: 2006--2026.
\end{tablenotes}
\end{threeparttable}
\end{table}

\end{document}
